\clearpage\fancyhead[L]{\small BORDA’S KREIS.}\hypertarget{native-409-part-lower}{}\label{native-409-part-lower}
\begin{center}\large BORDA’S KREIS.\end{center}
\begin{center}\hrule\vspace{4pt}\fontsize{10}{12}\selectfont V. \textsc{Liadenau} u. \textsc{Bohnenberger}, Zeitschrift für Astronomie u. verwandte Wiss.\\ Band V. Seite 198..211. März und April 1818.\par\vspace{4pt}\hrule\vspace{4pt}\end{center}
\par\indent Ihrem Wunsche zufolge, verehrtester Herr Professor, habe ich das Vergnügen, Ihnen hier diejenige Methode mitzutheilen, deren ich mich zur Parallelstellung der Gesichtslinie mit der Kreisebene bedient habe, und die, wie eine oft wiederholte Erfahrung mich gelehrt hat, einer sehr grossen Schärfe fähig ist.

\par\indent Um mich kurz und bestimmt ausdrücken zu können, bezeichne ich die verschiedenen Axen des Instruments auf folgende Weise. Diejenige Axe, die bei Höhenmessungen vertikal gestellt werden muss, und die man auch bei jedem andern Gebrauch wenigstens nicht viel von dieser Lage abweichen lässt, nenne ich die \emph{erste Axe}. Die Bewegung um diese Axe wird auf dem Azimuthkreise gemessen; da jedoch dieser nur Einen Index hat, dessen Vernier nur Minuten angibt, obwohl noch Theile von Minuten sich schätzen lassen, und man nicht sicher ist, ob nicht eine kleine Excentricität vorhanden ist, so kann eine Berichtigung der Gesichtslinie, vermittelst dieses Azimuthalkreises (indem man bei zwei entgegengesetzten vertikalen Lagen des Hauptkreises das Azimuth eines irdischen Gegenstandes, nöthigenfalls mit Rücksicht auf seine Parallaxe, wenn er nicht sehr entfernt ist, bestimmt) noch nicht alle erreichbare oder wünschenswerthe Genauigkeit geben; inzwischen mag man sich dieses bequemen und einfachen Verfahrens zur ersten vorläufigen Berichtigung bedienen.

\par\indent Nahe senkrecht zur ersten Axe ist diejenige, die ich die \emph{zweite} nenne; diese ist also jederzeit nahe horizontal. Durch Drehung um dieselbe wird der Kreis in die vertikale, schiefe oder horizontale Lage gebracht. An den \textsc{Reichenbach}schen Kreisen ist keine Theilung, um die Grösse dieser Drehung zu messen. Für die meisten Beobachtungen ist dies auch überflüssig; es gibt aber doch einige Fälle, wo es wünschenswerth wäre, namentlich würde man ohne Mühe Winkel zwischen Fixsternen und irdischen Objecten bei Tage messen können, was jetzt entweder gar nicht oder nur mit grosser Schwierigkeit zu bewerkstelligen ist.

\par\indent Wiederum fast senkrecht zur zweiten Axe ist die, um welche sich der ganze Kreis dreht, welche Bewegung am Trommelkreise gemessen wird. Da mit letzterer Axe, im mathematischen Sinn, diejenige, \SourceLeaf{410} um welche der Alhidadenkreis sich bewegt, zusammenfällt, so werden diese beiden Axen nur als Eine betrachtet, die ich die \emph{dritte} nenne. Diese dritte Axe ist auf der Kreisebene genau senkrecht, und wenn $a$ den Winkel der Gesichtslinie mit der Kreisebene bedeutet (welchen ich positiv setze, wenn die Objectivseite weiter von der Kreisfläche abliegt, als die Okularseite), so ist $90^\circ{-}a$ der Winkel der Gesichtslinie mit der dritten Axe, jene nach der Objectivseite, diese abwärts von der getheilten Fläche genommen.

\par\indent Für den Gebrauch des Instruments ist nicht nothwendig, dass die zweite Axe mit der Kreisebene genau parallel, und zur ersten Axe genau senkrecht sei; die Unterschiede werden jedoch gewiss klein sein, und eine oft wiederholte Erfahrung an dem hiesigen 12zolligen \textsc{Reichenbach}schen Kreise hat mich gelehrt, dass sie Jahre lang, und selbst, nachdem das Instrument auseinandergenommen gewesen, unveränderlich bleiben. Ich nenne $b$ die Neigung der zweiten Axe gegen die Kreisebene, positiv genommen, wenn der linke Arm\footnote{insofern man hinter dem senkrecht vorausgesetzten Kreise sich befindet.} von jener dem Kreise näher ist. Endlich nenne ich $c$ die Neigung der zweiten Axe gegen eine auf der ersten Axe senkrechten Ebene (wofür man die Ebene des Azimuthalkreises halten kann), positiv genommen, wenn der linke Arm mehr über dieselbe erhaben ist.

\par\indent Es seien noch $M$ und $N$ die Schnitte der dritten Axe mit der zweiten und mit der Gesichtslinie und $O$ ein gut zu sehendes Object. Die Entfernung $MN$ hat an jedem Kreise eine bestimmte, durch Abmessung oder auf andere Weise anzugebende Grösse, an dem hiesigen $\frac{1}{8}$ Meter. Insofern die Gesichtslinie auf das Object $O$ gerichtet ist, bezeichne ich den Winkel $MON$ mit $p$; die genaue Formel dafür ist

\[ \sin p = \frac{MN}{MO}\cdot\cos a \]

\noindent wofür man jederzeit $p = \dfrac{MN}{MO}\cdot 206265''$ annehmen darf. Es wird nicht schwer sein, $MO$ so genau zu erfahren, dass $p$ auf die Sekunde genau bekannt wird. Bei Objecten in bedeutender Entfernung wird $p$ ganz unmerklich; dies ist aber kein Grund, sie nähern bei der folgenden Berichtigungsmethode vorzuziehen, wenn letztere bessere Zielpuncte abgeben.

\par\indent Ich setze voraus, dass bei allen hier zu beschreibenden Operationen der Kreis selbst um die dritte Axe gar nicht bewegt wird, also die Hemmung an der Trommel fest angezogen ist, und die Stellschraube an derselben gar nicht berührt wird. Man kann, wenn man will, auch die Pressmutter am hintern Theile der dritten Axe etwas stärker als gewöhnlich anziehen.

\par\indent Ich nenne Kürze halber $I$ den Winkel zwischen einer Ebene durch die zweite und dritte Axe mit der Ebene durch die erste und zweite, welcher mithin die Grösse der Drehung des Kreises um die zweite Axe bestimmt. Der Bau des Instruments verstattet nur Werthe dieses Winkels von $0$ bis $90^\circ$, und etwas auf beiden Seiten über diese Grenzen hinaus; dass an den Reichenbachschen Kreisen keine Theilung zur Messung dieses Winkels sich befindet, habe ich schon oben bemerkt.

\par\indent Das Wesentliche meiner Methode beruht darauf, dass, wenn bei irgend einer Lage des Instruments die Gesichtslinie auf ein Object gerichtet ist, durch Verbindung einer Drehung um die zweite Axe mit einer Drehung des Alhidadenkreises um die dritte, die Gesichtslinie zum zweitenmale auf das Object geführt werden kann. Ich nehme an, es sei $I = i$, wenn die dritte Axe mit der zweiten und der geraden Linie $MO$ in Eine Ebene gebracht ist, der Index des Alhidadenkreises hingegen zeige $k$, wenn die Gesichtslinie mit der zweiten und der dritten Axe in Einer Ebene sich befindet, so dass die Gesichtslinie mit der zweiten Axe (linker Arm) auf der Ebene des Kreises einerlei Projection haben

\SourceLeaf{411}
\par\indent Man überzeugt sich leicht, dass wenn bei der Richtung der Gesichtslinie auf das Object der Index des Alhidadenkreises auf $k-\varkappa$ steht, und $I=i-\eta$ ist, die Gesichtslinie wiederum auf dieses Object gerichtet sein wird, wenn $I=i+\eta$ gemacht, und der Index des Alhidadenkreises auf $k+\varkappa$ gestellt wird; diese zweite Stellung wird natürlich nur dann physisch möglich sein, wenn der zweite Werth von $I$ innerhalb der Grenzen fällt, die die Einrichtung des Instruments zulässt. Es lässt sich ferner leicht darthun, dass nach aller Strenge folgende Gleichung Statt findet:

\[ \sin b \cos\varkappa = \cos b\, \mathrm{tang}(a+p) + \mathrm{cotang}\,\eta \sin\varkappa \]

\noindent und dass der Cosinus des Winkels der zweiten Axe mit der Linie $MO$

\[ = \sin b \sin(a+p) + \cos b \cos(a+p) \cos\varkappa \]

\noindent wird. Hieraus folgert man leicht, dass $\varkappa$ entweder sehr klein, oder sehr nahe $180^\circ$ sein wird; im ersten Fall wird der gedachte Winkel sehr klein, im zweiten nahe $180^\circ$ sein, d.\,i. im ersten Fall wird das Object nahe in der Richtung des linken Arms der zweiten Axe, im zweiten nahe in der Richtung des rechten liegen müssen.

\par\indent Man hat daher hinreichend genau

\noindent im ersten Fall \[ b - a - p = \frac{\varkappa}{\mathrm{tang}\,\eta} \]

\noindent im zweiten Fall \[ b + a + p = \frac{\varkappa - 180^\circ}{\mathrm{tang}\,\eta} \]

\par\indent Durch zwei solche Versuche erhält man daher, insofern man sich eine Kenntniss von $\eta$ verschaffen kann, die Werthe von $b-a-p$ und von $b+a+p$, und folglich von $b$ und $a+p$, oder, insofern $p$ als bekannt anzusehen ist, von $a$; zugleich erhält man offenbar aus jedem solchen Versuch (ohne $\eta$ nöthig zu haben) eine Bestimmung von $k$. Da $\eta$ hiebei willkührlich ist, so ist es am vortheilhaftesten, es so gross wie möglich zu machen. Indessen würde man doch nicht viel über $45^\circ$ hinaus gehen können, und da man an \textsc{Reichenbach}'schen Kreisen die Werthe von $I$ nicht ablesen kann, so muss man sich auf solche beschränken, wofür man die Stellung auf anderm Wege erhalten kann. Ich werde unten zeigen, wie man die Stellung, wofür $I=0$ ist, d.\,i. wo die erste. zweite und dritte Axe in Einer Ebene liegen, mit grösster Schärfe effectuiren kann. Man kann ferner die Kreisebene mit der ersten Axe parallel machen, indem man auf bekannte Weise beide vertikal stellt; man kann diese beiden Stellungen durch zwei auf dem Quadranten, an dem sich die Hemmung für die Drehung um die zweite Axe befindet, an der Kante der einen der Stützen, die die zweite Axe tragen, gezogene Striche bezeichnen, um sie immer schnell und hinreichend genau wieder finden zu können. Ich werde die erste Stelle Kürze halber die \emph{horizontale}, die zweite die \emph{vertikale} nennen. Für letztere ist nach aller Schärfe

\[ \cos I = -\,\mathrm{tang}\,b\,.\,\mathrm{tang}\,c \]

\noindent Es wird daher $I$ nicht merklich von $90^\circ$ verschieden sein, und da der Unterschied $=2\eta$ gesetzt ist, so wird genau

\[ \mathrm{tang}\,\eta = \BareRoot{\dfrac{1+\mathrm{tang}\,b\,.\,\mathrm{tang}\,c}{1-\mathrm{tang}\,b\,.\,\mathrm{tang}\,c}} = \BareRoot{\dfrac{\cos(b-c)}{\cos(b+c)}} \]

\noindent also nur unmerklich von $\eta$ verschieden sein.

\SourceLeaf{412}
\par\indent Hieraus folgt also, dass wenn bei dem ersten Versuch, wo das Object nahe in der Richtung des linken Arms der zweiten Axe sich befunden hat, bei der horizontalen Stellung des Kreises der Index des Alhidadenkreises $h$, bei der vertikalen $\upsilon$ gezeigt hat, im zweiten Versuch hingegen, wo das Object sich nahe in der Richtung des rechten Arms der zweiten Axe befand, der Index des Alhidadenkreises respectiv $180^\circ+h'$, $180^\circ+\upsilon'$ zeigte, man haben wird:

\[ \begin{aligned} b-a-p &= \upsilon-k = k-h = \tfrac{1}{2}(\upsilon-h) \\ b+a+p &= \upsilon'-k = k-h' = \tfrac{1}{2}(\upsilon'-h') \end{aligned} \]

\noindent folglich

\[ \begin{aligned} a &= \tfrac{1}{4}(\upsilon'-h'-\upsilon+h)-p \\ b &= \tfrac{1}{4}(\upsilon'-h'+\upsilon-h) \end{aligned} \]

\noindent Zugleich hat man für $k$ die doppelten Werthe

\[ \tfrac{1}{2}(\upsilon+h) \quad\text{und}\quad \tfrac{1}{2}(\upsilon'+h') \]

\noindent deren nahe Uebereinstimmung zur Bestätigung der Richtigkeit der Operationen dienen wird.

\par\indent Der Gang bei meinem Verfahren ist demnach folgender. Das Object, welches ich dabei wähle, muss einen scharfen Zielpunct darbieten und höchstens einige Grade vom Horizont entfernt sein. Ich stelle das Instrument so auf, dass der eine Fuss desselben mit der ersten Axe und dem Object nahe in Einer Ebene liegt, in welche ich sodann auch nach dem Augenmaass durch Azimuthalbewegung die zweite Axe bringe, sodass zuerst deren linker Arm dem Object zugekehrt ist. Hiernächst stelle ich den Kreis horizontal, und bringe durch die Schraube an dem erwähnten Fusse und durch Alhidadenbewegung die Gesichtslinie auf das Object. Sodann wird der Kreis vertikal gestellt, und die Gesichtslinie durch Azimuthal- und Alhidadenbewegung von neuem auf das Object gebracht. Hiedurch ist die erste grobe Stellung vollendet. Der Kreis wird von neuem horizontal gestellt, und falls die Alhidadenbewegung nicht hinreicht, um die Gesichtslinie auf das Object zu führen, an der Fussschraube nachgeholfen. Dann wird der Kreis wiederum vertikal gestellt, und die Gesichtslinie durch Alhidadenbewegung, nöthigenfalls mit Nachhülfe an der Azimuthalschraube, auf das Object gebracht. Mit diesen beiden Operationen wird abwechselnd so lange fortgefahren, bis die blosse Alhidadenbewegung in horizontaler Lage zum scharfen Pointiren hinreicht. Die Ablesungen am Alhidadenkreise geben $h$ und $\upsilon$.

\par\indent Hiernächst wird nach der letzten Vertikalstellung die Azimuthalhemmung gelöst, der Kreis $180^\circ$ um die erste Axe bewegt, so dass jetzt der rechte Arm der zweiten Axe nahe auf das Object gerichtet ist; durch Alhidadenbewegung (die nahe $180^\circ$ betragen wird) und feine Stellung an der Azimuthalschraube die Gesichtslinie scharf auf das Object gebracht; auch hier lese ich den Index des Alhidadenkreises ab, welcher $180^\circ+V$ geben mag. Dies ist der Anfang der zweiten Operationsreihe. Der Kreis wird horizontal gelegt, und durch Stellung an der Fussschraube und Alhidadenbewegung pointirt; bei der neuen Vertikalstellung wird Azimuthal- und Alhidadenbewegung angewandt. So wird wieder wechselsweise fortgefahren, bis die blosse Alhidadenbewegung zureicht, das Object in beiden Lagen zu treffen. Die Ablesungen am Alhidadenkreise geben hier $180^\circ+h'$ und $180^\circ+\upsilon'$. Die Messungen sind hiermit vollendet.

\par\indent Die vorhin angeführte Ablesung $180^\circ+V$ dient noch dazu, auch die Grösse $c$ zu bestimmen. Offenbar wird, bei vertikaler Stellung des Kreises, die Gesichtslinie mit der ersten und dritten Axe in Einer Ebene, und zwar nach oben zu gerichtet, sein, wenn der Index des Alhidadenkreises auf $\tfrac{1}{2}(\upsilon+V)-90^\circ$ steht, oder die Projection der Gesichtslinie auf der Kreisebene wird dann mit der ersten Axe parallel sein. Die Projection der zweiten Axe (linker Arm) auf eben dieser Ebene macht also mit dem Radius, der der ersten Axe parallel ist, den Winkel $k+90^\circ-\tfrac{1}{2}(\upsilon+V)$. Es ist daher

\SourceLeaf{413}
\[ \sin c = \cos b\,\sin\bigl(\tfrac{1}{2}(\upsilon+V)-k\bigr) \]

\noindent wofür man

\[ c = \tfrac{1}{2}(\upsilon+V)-k = \tfrac{1}{2}(V-h) \]

\noindent setzen darf.

\par\indent Dieses Verfahren führt also zur Kenntniss des Fehlers $a$ in Sekunden. Will man ihn wegschaffen, so muss das Fadenkreuz, insofern $a$ positiv ist, abwärts vom Kreise geschoben werden, oder, wenn man sich des prismatischen Okulars bedient, in der Richtung vom Objectiv zum Okular. Um die Grösse der Verschiebung anschaulich vor sich zu haben, kann man sich entweder einer Schätzung des Verhältnisses zu einem schicklichen Object von bekanntem Durchmesser nach dem Augenmaass bedienen, oder zuerst die Gesichtslinie scharf auf ein feines Object bringen, und dann den Alhidadenkreis um die Grösse $a$ verschieben, wo dann das Fadenkreuz soviel verrückt werden muss, bis das Object in der Diagonale erscheint. Dass man bei dieser delikaten Operation vorsichtig zu Werke gehen und alle Verrückung des Kreises selbst, während derselben, vermeiden muss, versteht sich von selbst. Immer bietet aber die Wiederholung der vorstehenden Methode das Mittel dar, sich zu versichern, ob man seinen Zweck erreicht habe.

\par\indent Es bleibt mir jetzt noch übrig, zu zeigen, wie die horizontale Stellung des Kreises, wo $I=0$, zu effectuiren sei. Insofern $c$ und $b$ nicht beträchtlich sind, kommt es freilich so genau nicht darauf an, und man könnte sich mit dem Augenmaasse begnügen. Allein durch folgendes Verfahren lässt sich dieses mit äusserster Schärfe erreichen.

\par\indent Nachdem bei festem Stande der Trommel $k$ bestimmt ist, lege man zuerst nach dem Augenmaass den Kreis horizontal, stelle den Index des Alhidadenkreises auf $k+l$, wo $l$ einen beliebigen Winkel, am besten $90^\circ$, bedeutet, und bringe durch Azimuthal- und Fussbewegung ein gut zu sehendes Object auf die Gesichtslinie. Hierauf stelle man den Index des Alhidadenkreises auf $k-l$; hat man die Horizontalstellung zufällig genau getroffen, so wird durch blosse Azimuthalbewegung die Gesichtslinie wieder genau auf das Object geführt werden können, sonst aber darüber oder darunter weggehen. Nachdem man dann durch die Azimuthalbewegung dies Object der Gesichtslinie so nahe wie möglich (d. i. auf den der dritten Axe parallelen, oder bei prismatischem Okular auf den mit dem Fernrohr parallelen Faden) geführt hat, bringe man durch Drehung des Kreises um die zweite Axe das Object der Gesichtslinie um die Hälfte näher und corrigire die andere Hälfte durch eine Fussschraube, nöthigenfalls mit einer Nachhülfe der Azimuthalschraube, so dass das Object wieder scharf auf der Gesichtslinie steht. Um die Genauigkeit der Halbirung zu prüfen, stelle man von neuem den Index des Alhidadenkreises auf $k+l$ und versuche, ob jetzt die Azimuthalbewegung allein das Object genau auf die Gesichtslinie zurückführt, wo nicht, so corrigirt man wieder die Hälfte des jetzt gewiss viel kleinern Fehlers an der Schraube ohne Ende, wodurch die feine Bewegung um die zweite Axe geschieht. Mit dieser Operation fährt man so lange fort, bis man vollkommen befriedigt ist.

\par\indent Wäre die erste, nach dem Augenmaasse gemachte, Stellung gar zu schlecht gerathen, so würde es schwer sein, die erste Halbirung leidlich genau zu machen, weil die Azimuthalbewegung, insofern $l$ gross angenommen wurde, das Object gar nicht in das Gesichtsfeld bringen würde. Um sich also vieles Probiren zu ersparen, nehme man Anfangs $l$ lieber klein, etwa $=10^\circ$, und bringe den Kreis nach vorgeschriebener Methode in die verlangte Stellung, welche denn durch Wiederholung des Verfahrens bei einem grossen Werth von $l$ fein berichtigt werden kann.

\par\indent Wenn an einem Instrument zufällig $b=c$ ist, so wird die Kreisebene bei horizontaler Stellung genau auf die erste Axe senkrecht sein, und ein Object, welches einmal auf der Gesichtslinie ist, wird bei jeder veränderten Alhidadenstellung durch blosse Azimuthalbewegung auf dieselbe zurückgebracht. \SourceLeaf{414} Sind aber $b$ und $c$ ungleich, so überzeugt man sich leicht, dass dies unmöglich ist; die Kreisebene kann dann im \emph{eigentlichen} Sinne nicht genau horizontal gestellt werden, wenn nicht die erste Axe etwas geneigt wird. Bei dem hiesigen \textsc{Reichenbach}schen Kreise ist $c=+5'40''$, $b=+28''$ im Mittel aus mehrern Versuchen.

\par\indent Uebrigens ist schon oben bemerkt, dass die Bestimmung von $k$ von der Kenntniss bestimmter Werthe von $I$ ganz unabhängig ist. Man braucht bloss, ohne Azimuthal- und Trommelschraube zu berühren, durch Verbindung von Alhidadenbewegung und Drehung um die zweite Axe ein Object zweimal auf die Gesichtslinie zu führen und das Mittel aus beiden Stellungen des Index des Alhidadenkreises zu nehmen. Am zweckmässigsten ist es hiebei, es so einzurichten, dass die Drehung um die zweite Axe so gross wie möglich wird. Um dies zu erhalten, mache ich vorläufig die oben gelehrte grobe Stellung, und wenn dann, bei ungefähr horizontaler Stellung des Kreises, das Object genau auf der Gesichtslinie ist, drehe ich den Kreis um die zweite Axe, bis die Bogenbewegung, die das Object hiedurch im Gesichtsfelde erhält, dasselbe wieder auf den Faden führt, der mit der Kreisebene parallel ist (bei prismatischem Okular auf den, welcher diesen repräsentirt, d. i. senkrecht zum Fernrohr steht), und bringe durch Alhidadenbewegung dann das Object genau auf das Kreuz. \emph{Bliebe} bei der Drehung um die zweite Axe das Object genau auf dem Kreuz, so wäre dies ein Beweis, dass das Object in der zweiten Axe selbst läge, dass folglich $b=a+p$, und die Stellung des Index selbst $=k$.

\par\indent Ich will bei dieser Gelegenheit noch des Verfahrens erwähnen, dessen ich mich mit Vortheil bediene, um die Kreisebene genau vertikal zu stellen. Bei \textsc{Reichenbach}s Kreisen ist dazu eigentlich die Hänglibelle bestimmt, die zwar grosse Genauigkeit geben kann, deren Gebrauch aber, weil das Fernrohr abgenommen werden muss, etwas umständlich ist, eine sehr solide Aufstellung und sehr viel Behutsamkeit erfordert. Ich bediene mich des Quecksilberhorizonts, wähle ein Object, so hoch, wie möglich (es braucht übrigens nicht entfernt zu sein, da es bei dieser Operation nichts schadet, wenn man auch die Okularröhre, um es deutlich zu sehen, weit herausziehen muss), bringe den nach dem Augenmaass vertikal gestellten Kreis in dessen Vertikalebene, und pointire genau auf dasselbe. Hierauf wird das Fernrohr heruntergeführt, bis man das Bild des Objects im Quecksilberhorizont sieht. Kommt die Gesichtslinie genau auf dasselbe, so steht die Kreisebene vertikal; wo nicht, so corrigirt man die eine Hälfte entweder an einer Fussschraube, oder an der Schraube ohne Ende, die andere an der Azimuthalschraube. Hat man die Halbirung genau getroffen; so wird die Gesichtslinie des wieder hinaufgeführten Fernrohrs das Object selbst genau treffen, sonst wird abermals die eine Hälfte der Abweichung an der Fussschraube oder der Schraube ohne Ende, die andere an der Azimuthalschraube corrigirt. Ist die Elevation des Objects nicht zu klein, und dieses scharf genug zu pointiren, so lässt sich auf diesem Wege auch eine sehr grosse Genauigkeit erreichen, und eine viel grössere, als nöthig ist, wenn man nicht \emph{sehr} kleine Zenithdistanzen messen will. Auf \emph{gar zu kleine} Zenithdistanzen ist es aber bekanntlich nicht rathsam, den \textsc{Borda}schen Kreis anzuwenden.

